\begin{abstract}
Для стохастичного моделювання поширення епідемій у великих
контактних мережах часто бракує інформації про конкретні
контакти між людьми. Водночас можуть бути доступні
агреговані дані про переміщення населення та зв'язки між
містами чи іншими спільнотами. Внутрішні контактні мережі
спільнот можна наближено відтворювати за відомими
структурними характеристиками, проте конкретні контакти
між ними залишаються невідомими. У роботі досліджено
двомасштабний метод стохастичного моделювання (\TS{}),
який поєднує детальні локальні симулятори Gillespie
з агрегованими міжспільнотними інтенсивностями
передавання інфекції без явного представлення
міжспільнотних контактів.

Теоретично доведено, що ідеальний двомасштабний процес
еквівалентний точному моделюванню на сурогатній мережі
з рівномірно розподіленими міжспільнотними зв'язками.
Визначено механізми розбіжностей із повномережевою
моделлю та встановлено стохастично не пізніший
час першого експорту інфекції за рівномірного
випадкового розміщення вузлів-мостів у розглянутій
теоретичній постановці.

Точність методу досліджено шляхом порівняння
з повномережевим Gillespie на синтетичних SI-мережах,
зокрема у 48 конфігураціях ланцюгів із десяти
спільнот. Експерименти продемонстрували близькість
часових розподілів за швидкого внутрішнього поширення:
у ланцюгу десяти клік нормована відстань Вассерштайна
після дев'яти міжспільнотних переходів становила
$3.34\%$. Встановлено залежність розбіжностей
від структури й розміру спільнот, кількості
міжспільнотних ребер та співвідношення часових
масштабів внутрішнього й міжспільнотного поширення.
Запропонований показник цього співвідношення
продемонстрував сильний ранговий зв'язок із величиною
розбіжностей (кореляція Спірмена $0.91$).

Результати підтверджують можливість відтворення
близьких до повномережевого моделювання часових
характеристик без знання конкретної структури
міжспільнотних контактів та допомагають визначати
умови, за яких таке наближення є доцільним.
\end{abstract}

\begin{keyword}
Стохастичне моделювання епідемій
\sep Алгоритм Gillespie
\sep Багатомасштабне моделювання
\sep Мережі спільнот
\sep Агреговані міжспільнотні зв'язки
\sep Неповні контактні дані
\sep Точність моделювання
\end{keyword}
